%Daten zur Institution

%\input{Dozentdaten}

%\renewcommand{\fachbereich}{Fachbereich}

%\renewcommand{\dozent}{Prof. Dr. . }

%Klausurdaten

\renewcommand{\klausurgebiet}{ }

\renewcommand{\klausurtyp}{ }

\renewcommand{\klausurdatum}{ . 20}

\klausurvorspann {\fachbereich} {\klausurdatum} {\dozent} {\klausurgebiet} {\klausurtyp}

%Daten für folgende Punktetabelle

\renewcommand{\aeins}{  3  }

\renewcommand{\azwei}{  3   }

\renewcommand{\adrei}{  3  }

\renewcommand{\avier}{  0  }

\renewcommand{\afuenf}{  2  }

\renewcommand{\asechs}{  0  }

\renewcommand{\asieben}{  3  }

\renewcommand{\aacht}{  7  }

\renewcommand{\aneun}{  0  }

\renewcommand{\azehn}{  7  }

\renewcommand{\aelf}{  4  }

\renewcommand{\azwoelf}{  3  }

\renewcommand{\adreizehn}{  3   }

\renewcommand{\avierzehn}{  7  }

\renewcommand{\afuenfzehn}{  2  }

\renewcommand{\asechzehn}{  9  }

\renewcommand{\asiebzehn}{  7  }

\renewcommand{\aachtzehn}{  63  }

\renewcommand{\aneunzehn}{    }

\renewcommand{\azwanzig}{    }

\renewcommand{\aeinundzwanzig}{    }

\renewcommand{\azweiundzwanzig}{    }

\renewcommand{\adreiundzwanzig}{    }

\renewcommand{\avierundzwanzig}{    }

\renewcommand{\afuenfundzwanzig}{    }

\renewcommand{\asechsundzwanzig}{    }

\punktetabellesiebzehn

\klausurnote

\newpage

\setcounter{section}{0}

\inputaufgabegibtloesung
{3}
{

Definiere die folgenden
\zusatzklammer {kursiv gedruckten} {} {} Begriffe.
\aufzaehlungsechs{Die
\stichwort {Gauß-Abbildung} {}
zu einer differenzierbaren Hyperfläche

\mavergleichskette
{\vergleichskette
{ Y
}
{ \subseteq }{ \R^n
}
{  }{
}
{    }{
}
{   }{
}
}
{}{}{.}

}{Eine
\stichwort {topologische Karte} {}
auf einer
\definitionsverweis {topologischen Mannigfaltigkeit}{}{}
$M$.

}{Eine  $C^1$-\stichwort {differenzierbare Mannigfaltigkeit} {} $M$.

}{Die $n$-te \stichwort {äußere Potenz} {} zu einem
$K$-\definitionsverweis {Vektorraum}{}{}
$V$
\zusatzklammer {es genügt, das Symbol dafür anzugeben} {} {.}

}{Eine
\stichwort {lokale Isometrie} {}
\maabb {\varphi} {L} {M
} {}
zwischen riemannschen Mannigfaltigkeiten.

}{Ein
\stichwort {lokal integrabler} {}
Zusammenhang auf einem differenzierbaren Vektorbündel
\maabb {p} {E} {X
} {}
über einer differenzierbaren Mannigfaltigkeit $X$.
}

}
{} {}

\inputaufgabegibtloesung
{3}
{

Formuliere die folgenden Sätze.
\aufzaehlungdrei{Der
\stichwort {Isometriesatz für den Paralleltransport auf einer Hyperfläche} {}

\mavergleichskette
{\vergleichskette
{  Y
}
{ \subseteq }{  \R^n
}
{  }{
}
{    }{
}
{   }{
}
}
{}{}{.}}{Das
\stichwort {Theorema egregium} {.}}{Der Satz über die Partition der Eins.}

}
{} {}

\inputaufgabegibtloesung
{3}
{

Es seien
\maabbdisp {f,g} {\R} { \R^n
} {}
\definitionsverweis {differenzierbare Kurven}{}{.}
Berechne die
\definitionsverweis {Ableitung}{}{}
der Funktion
\maabbeledisp {} {\R} {\R
} { t } { \left\langle f(t) , g(t)  \right\rangle
} {.}
Formuliere das Ergebnis mit dem Skalarprodukt.

}
{} {}

\inputaufgabe
{0}
{

}
{} {}

\inputaufgabegibtloesung
{2}
{

Es sei

\mavergleichskette
{\vergleichskette
{ Y
}
{ \subseteq }{ W
}
{ \subseteq }{ \R^n
}
{    }{
}
{   }{
}
}
{}{}{,}
$W$
\definitionsverweis {offen}{}{,}
eine
\definitionsverweis {differenzierbare Hyperfläche}{}{}
und sei
\maabb {\gamma} {I =[a,b]} {Y
} {}
eine
\definitionsverweis {geodätische Kurve}{}{.}
Zeige, dass der
\definitionsverweis {Paralleltransport}{}{}
längs $\gamma$ den tangentialen Vektor

\mavergleichskette
{\vergleichskette
{ \gamma'(a)
}
{ \in }{ T_{\gamma(a)}Y
}
{  }{
}
{    }{
}
{   }{
}
}
{}{}{}
in den tangentialen Vektor

\mavergleichskette
{\vergleichskette
{ \gamma'(b)
}
{ \in }{ T_{\gamma(b)}Y
}
{  }{
}
{    }{
}
{   }{
}
}
{}{}{}
überführt.

}
{} {}

\inputaufgabe
{0}
{

}
{} {}

\inputaufgabegibtloesung
{3}
{

Ist die Abbildung
\maabbeledisp {\varphi} {S^1} {S^1
} {(x,y)} {( \betrag {  x  } ,y)
} {,}
\definitionsverweis {differenzierbar}{}{?}

}
{} {}

\inputaufgabegibtloesung
{7}
{

Es sei $M$ eine kompakte differenzierbare Mannigfaltigkeit mit einer stetigen positiven Volumenform $\omega$. Zeige, dass

\mavergleichskettedisp
{\vergleichskette
{ \int_M \omega
}
{ <} { \infty
}
{ } {
}
{ } {
}
{ } {
}
}
{}{}{}
ist.

}
{} {}

\inputaufgabe
{0}
{

}
{} {}

\inputaufgabegibtloesung
{7 (1+2+4)}
{

\aufzaehlungdreiabc{ Man gebe eine topologische Eigenschaft an, die zeigt, dass das
\definitionsverweis {offene Einheitsintervall}{}{}
und die
\definitionsverweis {Kreissphäre}{}{}
$S^1$ nicht
\definitionsverweis {homöomorph}{}{}
sind.
}{Zeige, dass jede stetige
$1$-\definitionsverweis {Differentialform}{}{}
auf
\mathl{]0,1[}{}
\definitionsverweis {exakt}{}{}
ist.
}{ Man gebe eine konkrete
\definitionsverweis {geschlossene}{}{}
$1$-Form auf $S^1$ an, die nicht exakt ist.
}

}
{} {}

\inputaufgabegibtloesung
{4}
{

Es sei

\mavergleichskette
{\vergleichskette
{ U
}
{ \subseteq }{ \R^n
}
{  }{
}
{    }{
}
{   }{
}
}
{}{}{}
\definitionsverweis {offen}{}{}
und
\maabb {f} { U } { \R
} {}
eine zweifach
\definitionsverweis {stetig differenzierbare Funktion}{}{.}
Zeige

\mavergleichskettedisp
{\vergleichskette
{ d(df)
}
{ =} { 0
}
{ } {
}
{ } {
}
{ } {
}
}
{}{}{,}
wobei $d$ die
\definitionsverweis {äußere Ableitung}{}{}
bezeichnet.

}
{} {}

\inputaufgabegibtloesung
{3}
{

Beschreibe die
\definitionsverweis {Abbildung}{}{}
\maabbeledisp {\psi} { \Complex \setminus \{1\} } { \Complex
} {w} { { \frac{   (w+1)  { \mathrm i}  }{ 1-w  } }
} {}
in reellen Koordinaten.

}
{} {}

\inputaufgabegibtloesung
{3}
{

Es sei

\mavergleichskette
{\vergleichskette
{ H
}
{ \subseteq }{ \R^n
}
{  }{
}
{    }{
}
{   }{
}
}
{}{}{}
ein
\definitionsverweis {euklidischer Halbraum}{}{}
und

\mavergleichskette
{\vergleichskette
{ P
}
{ \in }{ H
}
{  }{
}
{    }{
}
{   }{
}
}
{}{}{.}
Es gebe eine in $\R^n$ offene Menge $V$ mit

\mavergleichskette
{\vergleichskette
{ P
}
{ \in }{ V
}
{ \subseteq }{ H
}
{    }{
}
{   }{
}
}
{}{}{.}
Zeige, dass
\mathl{P}{} kein Randpunkt von $H$ ist.

}
{} {}

\inputaufgabegibtloesung
{7}
{

Es sei $M$ eine $n$-dimensionale
\definitionsverweis {differenzierbare Mannigfaltigkeit}{}{}
und

\mavergleichskette
{\vergleichskette
{ P
}
{ \in }{ M
}
{  }{
}
{    }{
}
{   }{
}
}
{}{}{}
ein Punkt. Zeige, dass es eine
\definitionsverweis {Mannigfaltigkeit mit Rand}{}{}
$N$ gibt, deren Rand $\partial N$  diffeomorph zur $(n-1)$-dimensionalen Sphäre $S^{n-1}$ ist und derart, dass
\mathkor {} {M \setminus \{ P \}} {und} {N \setminus \partial N} {}
zueinander diffeomorph sind.

}
{} {}

\inputaufgabegibtloesung
{2}
{

Man gebe eine
\definitionsverweis {kompakte Ausschöpfung}{}{}
für den $\R^d$ an.

}
{} {}

\inputaufgabegibtloesung
{9}
{

Beweise den Brouwerschen Fixpunktsatz.

}
{} {}

\inputaufgabegibtloesung
{7 (1+2+4)}
{

Wir betrachten auf dem
\definitionsverweis {trivialen Vektorbündel}{}{}

\mathl{\R \times \R^2}{} über $\R$ den
\definitionsverweis {trivialen Zusammenhang}{}{}
$\nabla$.
\aufzaehlungdrei{Zeige, dass

\mavergleichskette
{\vergleichskette
{ f_1(t)
}
{ = }{ \left(  1+ t^2  , \,  t^3                   \right)
}
{  }{
}
{    }{
}
{   }{
}
}
{}{}{}
und

\mavergleichskette
{\vergleichskette
{ f_2(t)
}
{ = }{ \left(  t  , \,   1+t^2                   \right)
}
{  }{
}
{    }{
}
{   }{
}
}
{}{}{}
\definitionsverweis {Basisschnitte}{}{}
des Vektorbündels sind.
}{Drücke die Standardbasisschnitte $e_1,e_2$ als Linearkombination der Basisschnitte $f_1,f_2$ aus.
}{Bestimme die
\definitionsverweis {Christoffelsymbole}{}{}
von $\nabla$ bezüglich dieser Basisschnitte.
}

}
{} {}

 [[Kategorie:Latexseite]]
